\documentclass[12pt,CJK]{article}
\usepackage{geometry}
\usepackage{CJKutf8}
\usepackage{amsmath}
\usepackage{amssymb}
\usepackage{graphicx}
\usepackage{eufrak}
\usepackage{color}
\usepackage{slashed}
\usepackage{simplewick}
\usepackage{tikz}
\usepackage{tcolorbox}
\graphicspath{{../figures/}}
%%figures
\def\lfig#1#2{\includegraphics[width=#1 in]{#2}}
\def\addfig#1#2{\begin{center}\includegraphics[width=#1 in]{#2}\end{center}}
\def\wulian{\includegraphics[width=0.18in]{emoji_wulian.jpg}}
\def\bigwulian{\includegraphics[width=0.35in]{emoji_wulian.jpg}}
\def\bye{\includegraphics[width=0.18in]{emoji_bye.jpg}}
\def\bigbye{\includegraphics[width=0.35in]{emoji_bye.jpg}}
\def\huaixiao{\includegraphics[width=0.18in]{emoji_huaixiao.jpg}}
\def\bighuaixiao{\includegraphics[width=0.35in]{emoji_huaixiao.jpg}}
\def\jianxiao{\includegraphics[width=0.18in]{emoji_jianxiao.jpg}}
\def\bigjianxiao{\includegraphics[width=0.35in]{emoji_jianxiao.jpg}}
%% colors
\def\blacktext#1{{\color{black}#1}}
\def\bluetext#1{{\color{blue}#1}}
\def\redtext#1{{\color{red}#1}}
\def\darkbluetext#1{{\color[rgb]{0,0.2,0.6}#1}}
\def\skybluetext#1{{\color[rgb]{0.2,0.7,1.}#1}}
\def\cyantext#1{{\color[rgb]{0.,0.5,0.5}#1}}
\def\greentext#1{{\color[rgb]{0,0.7,0.1}#1}}
\def\darkgray{\color[rgb]{0.2,0.2,0.2}}
\def\lightgray{\color[rgb]{0.6,0.6,0.6}}
\def\gray{\color[rgb]{0.4,0.4,0.4}}
\def\blue{\color{blue}}
\def\red{\color{red}}
\def\green{\color{green}}
\def\darkgreen{\color[rgb]{0,0.4,0.1}}
\def\darkblue{\color[rgb]{0,0.2,0.6}}
\def\skyblue{\color[rgb]{0.2,0.7,1.}}
%%control
\def\be{\begin{equation}}
\def\ee{\nonumber\end{equation}}
\def\bea{\begin{eqnarray}}
\def\eea{\nonumber\end{eqnarray}}
\def\bch{\begin{CJK}{UTF8}{gbsn}}
\def\ech{\end{CJK}}
\def\bitem{\begin{itemize}}
\def\eitem{\end{itemize}}
\def\bcenter{\begin{center}}
\def\ecenter{\end{center}}
\def\bex{\begin{minipage}{0.2\textwidth}\includegraphics[width=0.6in]{jugelizi.png}\end{minipage}\begin{minipage}{0.76\textwidth}}
\def\eex{\end{minipage}}
\def\chtitle#1{\frametitle{\bch#1\ech}}
\def\bmat#1{\left(\begin{array}{#1}}
\def\emat{\end{array}\right)}
\def\bcase#1{\left\{\begin{array}{#1}}
\def\ecase{\end{array}\right.}
\def\bmini#1{\begin{minipage}{#1\textwidth}}
\def\emini{\end{minipage}}
\def\tbox#1{\begin{tcolorbox}#1\end{tcolorbox}}
\def\pfrac#1#2#3{\left(\frac{\partial #1}{\partial #2}\right)_{#3}}
%%symbols
\def\bropt{\,(\ \ \ )}
\def\optlist#1#2#3{(A)\;#1\ \ \ (B)\;#2\ \ \ (C)\;#3}
\def\sone{$\star$}
\def\stwo{$\star\star$}
\def\sthree{$\star\star\star$}
\def\sfour{$\star\star\star\star$}
\def\sfive{$\star\star\star\star\star$}
\def\rint{{\int_\leftrightarrow}}
\def\roint{{\oint_\leftrightarrow}}
\def\stdHf{{\textit{\r H}_f}}
\def\deltaH{{\Delta \textit{\r H}}}
\def\ii{{\dot{\imath}}}
\def\skipline{{\vskip0.1in}}
\def\skiplines{{\vskip0.2in}}
\def\lagr{{\mathcal{L}}}
\def\hamil{{\mathcal{H}}}
\def\vecv{{\mathbf{v}}}
\def\vecx{{\mathbf{x}}}
\def\vecy{{\mathbf{y}}}
\def\veck{{\mathbf{k}}}
\def\vecp{{\mathbf{p}}}
\def\vecn{{\mathbf{n}}}
\def\vecA{{\mathbf{A}}}
\def\vecP{{\mathbf{P}}}
\def\vecsigma{{\mathbf{\sigma}}}
\def\hatJn{{\hat{J_\vecn}}}
\def\hatJx{{\hat{J_x}}}
\def\hatJy{{\hat{J_y}}}
\def\hatJz{{\hat{J_z}}}
\def\hatj#1{\hat{J_{#1}}}
\def\hatphi{{\hat{\phi}}}
\def\hatq{{\hat{q}}}
\def\hatpi{{\hat{\pi}}}
\def\vel{\upsilon}
\def\Dint{{\mathcal{D}}}
\def\adag{{\hat{a}^\dagger}}
\def\bdag{{\hat{b}^\dagger}}
\def\cdag{{\hat{c}^\dagger}}
\def\ddag{{\hat{d}^\dagger}}
\def\hata{{\hat{a}}}
\def\hatb{{\hat{b}}}
\def\hatc{{\hat{c}}}
\def\hatd{{\hat{d}}}
\def\hatN{{\hat{N}}}
\def\hatH{{\hat{H}}}
\def\hatp{{\hat{p}}}
\def\Fup{{F^{\mu\nu}}}
\def\Fdown{{F_{\mu\nu}}}
\def\newl{\nonumber \\}
\def\vece{\mathrm{e}}
\def\calM{{\mathcal{M}}}
\def\calT{{\mathcal{T}}}
\def\calR{{\mathcal{R}}}
\def\barpsi{\bar{\psi}}
\def\baru{\bar{u}}
\def\barv{\bar{\upsilon}}
\def\qeq{\stackrel{?}{=}}
\def\torder#1{\mathcal{T}\left(#1\right)}
\def\rorder#1{\mathcal{R}\left(#1\right)}
\def\contr#1#2{\contraction{}{#1}{}{#2}#1#2}
\def\trof#1{\mathrm{Tr}\left(#1\right)}
\def\trace{\mathrm{Tr}}
\def\comm#1{\ \ \ \left(\mathrm{used}\ #1\right)}
\def\tcomm#1{\ \ \ (\text{#1})}
\def\slp{\slashed{p}}
\def\slk{\slashed{k}}
\def\calp{{\mathfrak{p}}}
\def\veccalp{\mathbf{\mathfrak{p}}}
\def\Tthree{T_{\tiny \textcircled{3}}}
\def\pthree{p_{\tiny \textcircled{3}}}
\def\dbar{{\,\mathchar'26\mkern-12mu d}}
\def\erf{\mathrm{erf}}
\def\const{\mathrm{constant}}
\def\pheat{\pfrac p{\ln T}V}
\def\vheat{\pfrac V{\ln T}p}
%%units
\def\fdeg{{^\circ \mathrm{F}}}
\def\cdeg{^\circ \mathrm{C}}
\def\atm{\,\mathrm{atm}}
\def\angstrom{\,\text{\AA}}
\def\SIL{\,\mathrm{L}}
\def\SIkm{\,\mathrm{km}}
\def\SIyr{\,\mathrm{yr}}
\def\SIGyr{\,\mathrm{Gyr}}
\def\SIV{\,\mathrm{V}}
\def\SImV{\,\mathrm{mV}}
\def\SIeV{\,\mathrm{eV}}
\def\SIkeV{\,\mathrm{keV}}
\def\SIMeV{\,\mathrm{MeV}}
\def\SIGeV{\,\mathrm{GeV}}
\def\SIcal{\,\mathrm{cal}}
\def\SIkcal{\,\mathrm{kcal}}
\def\SImol{\,\mathrm{mol}}
\def\SIN{\,\mathrm{N}}
\def\SIHz{\,\mathrm{Hz}}
\def\SIm{\,\mathrm{m}}
\def\SIcm{\,\mathrm{cm}}
\def\SIfm{\,\mathrm{fm}}
\def\SImm{\,\mathrm{mm}}
\def\SInm{\,\mathrm{nm}}
\def\SImum{\,\mathrm{\mu m}}
\def\SIJ{\,\mathrm{J}}
\def\SIW{\,\mathrm{W}}
\def\SIkJ{\,\mathrm{kJ}}
\def\SIs{\,\mathrm{s}}
\def\SIkg{\,\mathrm{kg}}
\def\SIg{\,\mathrm{g}}
\def\SIK{\,\mathrm{K}}
\def\SImmHg{\,\mathrm{mmHg}}
\def\SIPa{\,\mathrm{Pa}}

\geometry{tmargin=0.3in, bmargin=0.5in, lmargin=0.5in, rmargin=0.9in, nohead, nofoot}
\def\mark#1{{\color{blue} (#1分)}}
\renewcommand{\thepage}{}
\begin{document}
\bch
{\large 数理方法 课堂小测II}

{\vskip 0.25in}

姓名 ....................... {\hskip 0.5in}    学号 .......................{\hskip 0.5in}  分数 ...................

{\vskip 0.1in}

\bitem
\item[(一)]{选择题，每题3分，共45分。

  \bitem

\item[(1)]{ $\ii^\ii$ 的所有可能取值为 \brans{D}

  \optlist{$-1$}{$\pm 1$}{$e^{-\frac{\pi}{2}}$}{$e^{-\left(2n+\frac{1}{2}\right)\pi},\ n\in Z$}
}

\item[(2)]{在幅角连续变化的意义下，$\sqrt{z}$ 的原函数为 (忽略不写积分常数) \brans{A}

  \optlist{$\frac{2}{3}z^{\frac{3}{2}}$}{$\frac{1}{2\sqrt{z}}$}{$\frac{4}{3}z^{\frac{3}{4}}$}{$\frac{1}{2}\ln z$}}
  

\item[(3)]{方程 $z^5 + z^4  + z^3+ 2=0$ 的所有复数根的倒数之和为 \brans{A}

  \optlist{$0$}{$1$}{$-1$}{$-2$}

  {\small \red 直接利用零次项和一次项的根和系数关系；或者令 $u=\frac{1}{z}$，分析$u$的方程的所有根的和.}
}


\item[(4)]{下列哪一个复数在单位圆 $|z|=1$ 的内部 \brans{A}

  \optlist{$\frac{5}{6} + \ii\cos 1 $}{$\frac{1}{2} + \ii \cos\frac{1}{2}$}{$\cos 2+ \ii \sin 2$}{$\frac{1}{2}-\ii$}

  {\small \red 利用 $\frac{1}{2}>\sin\frac{1}{2}$排除B；或利用 $\frac{5}{6} = 1-\frac{1}{6} < 1-\frac{1}{6}+\frac{1}{5!}-\frac{1}{7!}+\ldots = \sin 1$ 确定A正确}
  }
  
  
\item[(5)]{$z_1,z_2,z_3,z_4$ 是四个互不相同的复数，且 $\frac{(z_1-z_2)(z_3-z_4)}{(z_2-z_3)(z_4-z_1)}$ 为实数。那么 $z_1,z_2,z_3,z_4$ 在复平面上对应的四个点 \brans{D}

  \optlist{构成平行四边形}{共圆}{构成平行四边形或共线}{共圆或共线}

  {\small \red 画图，利用商的幅角的几何意义以及圆的几何性质。}
}

\item[(6)]{ $\frac{1}{(z+1)(z+2)}$ 在环形区域 $1<|z|<2$ 内Laurent展开的 $\frac{1}{z^2}$ 的系数为 \brans{C}

    \optlist{$0$}{$1$}{$-1$}{$-3$}}
  
\item[(7)]{$\frac{\sin^2 z}{\left(z-\pi\right)^5}$ 在 $z=\pi$ 处的留数等于 \brans{D}

  \optlist{$0$}{$\frac{2}{3}$}{$-\frac{2}{3}$}{$-\frac{1}{3}$}


  {\small \red 作变量替换 $t=z-\pi$并对分子进行展开；或者把$z=\pi$当成5阶极点用去极点的留数公式(虽然它实质上3阶极点，但课上讲过这样操作没有bug)。}
}


\item[(8)]{下列哪个多值函数在区域 $1<|z|<2$ 内可以取单值分枝成为解析函数 \brans{D}

    \optlist{$\ln\frac{z}{z-2}$}{$\sqrt{z(z-2)}$}{$\ln{\frac{z-1}{z-2}}$}{$\sqrt{z(z-1)(z-2)}$}}

\item[(9)]{ 考虑以原点为中心的一段小圆弧上的积分。当圆弧半径趋向于零时，下列哪个函数的积分一定趋向于零？  \brans{C}

  \optlist{$\frac{e^z}{z}$}{$\frac{1}{z^2}$}{$\frac{1}{\sqrt{z}}$的任意单值分枝}{$\frac{\ln z}{z}$的任意单值分枝}

{\small \red $\frac{1}{\sqrt{z}}$的原函数为 $2\sqrt{z}$，其变化量趋向于零。一般地，只有$z^\alpha$ ($\alpha>-1$)或者一个有界函数的小圆弧积分才能直接不加计算地忽略。}
}

\item[(10)]{$z=0$是下列哪一个函数的本性奇点（即邻域Laurent展开有无穷多个负次幂项)? \brans{D}

    \optlist{$\frac{z}{1-\cos z}$}{$\frac{1}{z^{1000}}$}{$\frac{1}{e^z-1}$}{$e^{\frac{1}{z}}$}

    {\red\small 选项D可以直接展开验证。A，C均满足$\lim_{z\rightarrow 0} z f(z)$有限，都是一阶极点。}
  }
  
\item[(11)]{ 记 $f(x)=e^{-\frac{x^4}{2}}$ 的傅立叶变换为 $F(k)$ ，则积分$\int_{-\infty}^{\infty} |F(k)|^2 dk$等于 \brans{A}

  \optlist{$\frac{1}{2}\sfgamma{\frac{1}{4}}$ }{$\frac{1}{2}\sfgamma{\frac{1}{3}}$}{$\frac{1}{2}\sqrt{\pi}$}{$\frac{1}{2}$}

{\small \red 利用傅立叶变换保持内积不变的性质$\int_{-\infty}^{\infty} |F(k)|^2 dk = \int_{-\infty}^\infty |f(x)|^2dx$.}
}
  
\item[(12)]{ 下列哪个数最大？ \brans{A}

  \optlist{$\sfgamma{\frac{1}{5}}$}{$\sfgamma{\frac{6}{5}}$}{$\sfgamma{\frac{11}{5}}$}{$\sfgamma{\frac{16}{5}}$}

  {\small \red 利用递推公式算出选项之间的比值。}
}

\item[(13)]{随机抛6000次骰子，恰好 1，2，3，4，5，6 每个面向上都是1000次的概率和下列哪个数量级最接近？ \brans{C}

  \optlist{$10^{-3}$}{$10^{-6}$}{$10^{-9}$}{$10^{-12}$}

{\small \red 对$P = \frac{\frac{6000!}{(1000!)^6}}{6^{6000}}$用Stirling公式.}
}

\item[(14)]{三维直角坐标系中，曲面$x^4+y^4+z^4=1$包围的体积为  \brans{D}

  \optlist{$\frac{\sfgamma{\frac{1}{4}}}{6\sqrt{2}\pi}$}{$\frac{\sfgamma{\frac{1}{4}}^2}{6\sqrt{2}\pi}$}{$\frac{\sfgamma{\frac{1}{4}}^3}{6\sqrt{2}\pi}$}{$\frac{\sfgamma{\frac{1}{4}}^4}{6\sqrt{2}\pi}$}

{\small \red 利用 $n$ 维限和积分公式以及互余宗量关系 $\sfgamma{\frac{1}{4}}\sfgamma{\frac{3}{4}} = \frac{\pi}{\sin\frac{\pi}{4}}$. }
}

\item[(15)]{已知全平面上解析的函数满足 $f(0) = 1$， $ |f(z)| > 2|z|-1 $。那么，下列哪个区域内一定有$f(z)$的零点？ \brans{A}

  \optlist{$|z|<1$}{$|z|>1$}{$|2z-1|<1$}{$|2z-1|>1$}

{\red\small 假设在单位圆内无零点，取解析函数 $g(z)=\frac{1}{f(z)}$，对$g(z)$运用最大模原理得到矛盾。}
}
  
  \eitem
  }
\item[(二)]{ 填空题，每题5分，共40分。


  \bitem
\item[(1)]{复变函数 $f(z)= ze^z$ 的原函数为 \ans{$(z-1)e^z+c$}.}      
\item[(2)]{近似到小数点后两位 $ \cos \frac{\ii}{10} \approx$ \ans{$1.01$} .}
\item[(3)]{逆时针方向沿着单位圆的围道积分$\oint_{|z|=1}\frac{1}{(10z^{6}-1)(z-2)}\,dz$等于 \ans{$-\frac{2\pi\ii}{639}$}.}
\item[(4)]{规定在 $z=1$ 处 $z$ 的幅角为零，且幅角连续变化。逆时针方向沿着上半个单位圆（从$1$到$-1$）的积分$\int_{|z|=1, \mathrm{Im}(z)\ge 0}\, \ln z \,dz = $ \ans{$2-\ii \pi$}.}
\item[(5)]{ 函数 $f(t) = \delta(t^2-1)$ 的拉普拉斯变换为 \ans{$\frac{e^{-p}}{2}$}.

{\small \red 注意拉普拉斯变换积分范围是$(0,\infty)$.}
}
\item[(6)]{记 $ F(k) = \int_{-\infty}^\infty \frac{\cos{\left[k(x^2-1)\right]}}{1+x^4} dx $，则 $ \int_{-\infty}^\infty F(k)\, dk = $ \ans{$\pi$} .

{\small \red 交换积分次序，并利用 $\int_{-\infty}^\infty e^{\ii k(x^2-1)}dk  = 2\pi \delta(x^2-1)$.}
}
\item[(7)]{$f(z) = \frac{z^8}{z^9+z^8+1}$的所有孤立奇点处的留数之和等于 \ans{$1$}.

{\red\small 在大圆上积分趋向于 $\oint \frac{dz}{z} = 2\pi \ii$，所以内部留数之和为1。}
}
\item[(8)]{实积分 $\int_0^{\frac{\pi}{2}} \ln (\sin x) \, dx  = $ \ans{$-\frac{\pi}{2}\ln 2$} .

  {\red\small 令$I=\int_0^{\frac{\pi}{2}} \ln(\sin x)dx$，根据对称性$I=\int_0^{\frac{\pi}{2}} \ln(\cos x)dx$。两式相加并利用倍角公式。}
}
  \eitem
}

\item[(三)]{函数$f(t)$满足微分方程
  $$f'' + 2f' + f =  2\cos t $$
  和初始条件
  $$f(0) = 1, f'(0) = -1. $$
  \bitem
\item[(1)]{设$f(t)$的拉普拉斯变换为$F(p)$，写出$F(p)$满足的代数方程，并解出$F(p)$； (10分)}
\item[(2)]{把$F(p)$逆变换求出$f(t)$。 (5分)


  \skipline

  {\blue
    \bitem
    \item[(1)]{
    $$(p^2F-p+1) + 2(pF-1) + F = \frac{2p}{p^2+1} ,$$
    解出
    $$F(p)=\frac{1}{p+1}+ \frac{2p}{(p^2+1)(p+1)^2}.$$}
    \item[(2)]{把$F(p)$拆写为
      $$F(p)=\frac{1}{p+1}+ \frac{1}{p^2+1}-\frac{1}{(p+1)^2}.$$
      逐项求逆变换，得到
      $$f(t) = e^{-t} + \sin t - te^{-t} $$

    }
      \eitem

    

    }

}  
  \eitem
}  
  
\eitem


\ech
\end{document}
